\section{Discussion and Conclusion}
\label{sec:discussion}
\label{sec:conclusion}

When every turn is observed, an input-dependent rotation can represent an exact
count, and a policy can learn it from a terminal reward. A learned rotation angle,
however, accumulates small errors over time, and a categorical counter avoids
them by fixing each shift and weighting the shifts by event probabilities. When
turns are missing, a cyclic filter keeps several possible positions and updates
their probabilities whenever a pointer reading arrives. In our tests the cyclic
filters gave the most accurate beliefs among the learned models. Because the same
planner often chose the same action from different beliefs, these better beliefs
changed net utility by only about one point. Their clearest effect was on whether
to take another reading: when readings were cheap and reliable, an overconfident
estimate could lead the planner to skip a useful one.

These comparisons are controlled in three ways. We supply the cycle size, and in
the feedback study we train estimators with state labels and give the planner the
simulator's actuator and sensor models. A programmed counter shows what is
possible when every event is known, and a Bayes filter on the observed history
shows how much uncertainty remains when the observation model is known. Giving
every estimate to the same planner then tests whether better beliefs change
actions or costs. On real sensors, learning a bias state and a disturbance-aware
pointer model would be the next step toward the classical filter.

The recorded inertial sensor shows the same pattern outside simulation. A rotation
set by the measured heading increment kept its accuracy from 10~s training
windows to 240~s test windows, whereas recurrent and attention models given the
same inputs did not. Under magnetic disturbance our learned pointer model only
matched these networks, and an engineered filter with bias estimation and
disturbance rejection remained more accurate in both tasks. Taken together, the
results suggest that a known cycle helps most when it fixes the entry point of
uncertainty: the transition when increments are measured, and the pointer model
when readings are corrupted.
